% -----------------------------------------------------------------------------
\section{Implementation, Tools, and Technologies}
\label{sec:implementation}
% -----------------------------------------------------------------------------

The final MoonAI codebase combines systems programming, GPU execution, UI rendering,
serialization, experiment analysis, and documentation tooling. The value of this section is not
to list technologies mechanically, but to explain why they were used.

\subsection{Core Technology Stack}

\begin{table}[H]
  \centering
  \caption{Main tools and technologies used in the final implementation}
  \begin{tabularx}{\textwidth}{@{}p{2.8cm}p{2.8cm}X@{}}
    \toprule
    Layer & Technology & Role in MoonAI \\
    \midrule
    Systems language & Rust 2024 & Primary implementation language for the runtime, UI,
    configuration, metrics, and analysis integration points \cite{rust}. \\
    GPU execution & CUDA via \texttt{build.rs} and \texttt{cc} & Executes the hot simulation,
    evolution, and readback path on NVIDIA GPUs \cite{cuda}. \\
    Desktop UI & \texttt{eframe} / \texttt{egui} & Provides the application shell, controls,
    panels, and interactive runtime workflow \cite{egui}. \\
    GPU rendering & \texttt{wgpu} & Renders the world view and supports custom population
    drawing in the UI \cite{wgpu}. \\
    Runtime configuration & \texttt{mlua} & Loads \texttt{experiments.lua} so simulation
    presets remain editable without recompilation \cite{mlua}. \\
    Persistence & \texttt{serde} / \texttt{serde\_json} & Stores settings and exports structured
    artifacts in JSON form \cite{serde}. \\
    Analysis & Pandas, Matplotlib, NetworkX, Jinja2 & Loads run artifacts, computes summaries,
    renders plots and genome diagrams, and builds self-contained HTML reports
    \cite{pandas,matplotlib,networkx,jinja2}. \\
    Workflow & \texttt{just}, \texttt{uv} & Standardizes build, test, and analysis commands for
    the project workflow. \\
    Documentation & Zensical, GitHub Pages, Mermaid, LaTeX & Publishes docs, renders diagrams,
    and produces course reports \cite{mermaid}. \\
    \bottomrule
  \end{tabularx}
\end{table}

\subsection{Why Rust and CUDA Were Used}

Rust was selected because the final system needed strong module organization, explicit data
types, good tooling, and direct interoperability with low-level code. The CUDA side was needed
because the project is fundamentally population-heavy: sensor computation, inference,
movement, reproduction, and compact readback become expensive when tens of thousands of agents
are simulated simultaneously.

The combination allows MoonAI to keep high-level application flow readable while still moving
the computationally dominant path to the GPU.

\subsection{Why the UI Stack Matters}

The UI stack is not cosmetic. It is part of the engineering solution. \texttt{egui} and
\texttt{wgpu} were chosen because the project required:

\begin{itemize}[nosep]
  \item a desktop application with live controls,
  \item custom rendering for a dense simulated world,
  \item selection and inspection of individual agents,
  \item in-application queue management and settings editing.
\end{itemize}

This turns MoonAI from a raw simulation executable into an actual experimental workbench.

\subsection{Why the Analysis Stack Matters}

The Python analysis stack was selected to support post-run comparison rather than real-time
simulation. Pandas is used for structured metric handling, Matplotlib for plots, NetworkX for
topology visualization, and Jinja2 for self-contained HTML report rendering. This separation of
concerns is pragmatic: Rust and CUDA handle the runtime, while Python handles data analysis and
presentation efficiently.

\subsection{Supporting Engineering Tools Used During the Course}

Several supporting tools were also important during the course:

\begin{itemize}[nosep]
  \item \textbf{\texttt{just}} standardized recurring commands such as build, test, and
    analysis.
  \item \textbf{\texttt{uv}} simplified Python environment and dependency management for the
    analysis tooling.
  \item \textbf{Mermaid CLI} was used to generate architecture diagrams for the final report and
    documentation \cite{mermaid}.
  \item \textbf{GitHub Pages} and Zensical were used to publish project documentation and report
    links.
  \item \textbf{LaTeX} was used to prepare the formal design and final course reports.
\end{itemize}

Together, these tools made the project easier to build, reason about, validate, and present.
